% Part: lambda-calculus
% Chapter: church-rosser
% Section: beta-eta-reduction

\documentclass[../../../include/open-logic-section]{subfiles}

\begin{document}

\olfileid{lam}{cr}{be}

\olsection{اختزال~$\beta\eta$}

ننتقل إلى إثبات خاصية تشيرش--روسر لاختزال~$\beta\eta$ ($\bered$).

\begin{lem}\ollabel{lem:one-par}
  من~$M \beredone M'$ يلزم~$M \beredpar M'$.
\end{lem}

\begin{proof}
  نستقرئ~!!{derivation} العبارة~$M \beredone M'$.
  فإن كان~$M \beredone M'$ حاصلًا بتحويل~$\eta$، أي بقاعدة
  \olref[syn][eta]{defn:beredone}، جمعنا
  \olref[pbe]{defn:beredpar5} إلى~\olref[pbe]{thm:refl}.
  وأما بقية الحالات فنجري فيها على برهان~\olref[b]{lem:one-par}.
\end{proof}


\begin{lem}\ollabel{lem:par-red}
  من~$M \beredpar M'$ يلزم~$M \bered M'$.
\end{lem}

\begin{proof}
  نستقرئ~!!{derivation} العبارة~$M \beredpar M'$.

  أما الحالات الأربع الأولى فبرهانها كما في~\olref[b]{lem:par-red}.
  وأما إذا كانت القاعدة الأخيرة~\olref[pbe]{defn:beredpar5}، فإن~$M$
  هو~$\lambd[x][Nx]$ و$M'$ هو~$N'$، لبعض~$x$ و$N$ و$N'$،
  مع~$x \notin \FV{N}$ و$N \beredpar N'$.
  نختزل حينئذ~$\lambd[x][Nx]$ إلى~$N$ بتحويل~$\eta$،
  ثم نجري خطوات~$\beredone$ التي تحقق~$N \bered N'$؛
  ووجود هذه المتتالية ثابت بفرض الاستقراء.
\end{proof}


\begin{lem}\ollabel{lem:str}
  العلاقة~$\bered$ أصغر علاقة متعدية تشتمل على~$\beredpar$.
\end{lem}

\begin{proof}
  يجرى البرهان على مثال~\olref[b]{lem:str}.
\end{proof}

\begin{thm}\ollabel{thm:cr}
  العلاقة~$\bered$ ذات خاصية تشيرش--روسر.
\end{thm}

\begin{proof}
  نضم~\olref[dap]{thm:str} و\olref[pbe]{thm:cr} إلى~\olref{lem:str}، فيتم المطلوب.
\end{proof}
\end{document}
